\section{Block calculus and reset equations}

This chapter develops the block calculus used by the unified core and
the corrected window bound.  It records the small-block formulas, the
reset equations, the first-block terminal algebra, and the block-head
identities.

\subsection{Block molecule formulas}

For a block with step weights $u_1,\dots,u_L$ and total weight
\begin{equation*}
  U=u_1+\cdots+u_L,
\end{equation*}
the block molecule is the word molecule of the block word.  The
recursive definition gives the following closed forms.

\begin{center}
\small
\begin{tabular}{@{}lll@{}}
\toprule
block word & molecule $B$ & total weight $U$ \\
\midrule
$[u_1]$ & $1$ & $u_1$ \\
$[u_1,u_2]$ & $5+2^{u_1}$ & $u_1+u_2$ \\
$[u_1,u_2,u_3]$ & $25+5\cdot2^{u_1}+2^{u_1+u_2}$ & $u_1+u_2+u_3$ \\
\bottomrule
\end{tabular}
\end{center}

The block equation is
\begin{equation*}
  2^U\,r_L=5^L\,r_0+B.
\end{equation*}
It is the exact word-orbit identity applied to the block word; no
approximation is involved.

\subsection{Derivation of the formulas}

For a singleton word, $\wordA([u_1])=1$ and $U=u_1$.  For length two,
\begin{equation*}
  \wordA([u_1,u_2])
    =2^0\cdot5^1+2^{u_1}\cdot5^0
    =5+2^{u_1}.
\end{equation*}
For length three,
\begin{equation*}
  \wordA([u_1,u_2,u_3])
    =2^0\cdot5^2+2^{u_1}\cdot5^1+2^{u_1+u_2}\cdot5^0
    =25+5\cdot2^{u_1}+2^{u_1+u_2}.
\end{equation*}
These are the exact forms substituted into the segment equations in
the unified-core chapter.

\subsection{Worked block examples}

The block starting at $7$ with weights $[2,1]$ has molecule $9$:
\begin{equation*}
  2^3\cdot23=184=5^2\cdot7+9.
\end{equation*}
The block starting at $9$ with weights $[1,2]$ has molecule $7$:
\begin{equation*}
  2^3\cdot29=232=5^2\cdot9+7.
\end{equation*}
The block starting at $29$ with weights $[1,1,2]$ has molecule
\begin{equation*}
  B=25+5\cdot2+2^{1+1}=39,
\end{equation*}
and
\begin{equation*}
  2^4\cdot229=3664=5^3\cdot29+39.
\end{equation*}

\subsection{Reset equations}

A reset step of weight $t\in\{1,2\}$ maps the candidate predecessor
$x$ to the block head $r_j$:
\begin{equation*}
  2^t\,r_j=5x+1.
\end{equation*}
The reset predecessor has the form
\begin{equation*}
  x=5^{k_0}\,s+\delta\,5^{j-1}-1,
\end{equation*}
where $s5^{k_0}=r+1$ for the previous terminal state $r$, and
$\delta=1$ for $t=1$, $\delta\in\{1,3\}$ for $t=2$.

Substituting into $2^tr_j=5x+1$ gives
\begin{equation*}
  2^t\,r_j=5^{k_0+1}s+\delta5^j-4.
\end{equation*}
Adding $2^t$ to both sides:
\begin{itemize}[leftmargin=*]
\item for $t=2$,
  \begin{equation*}
    4(r_j+1)=5^{k_0+1}s+\delta5^j;
  \end{equation*}
\item for $t=1$,
  \begin{equation*}
    2(r_j+1)=5^{k_0+1}s+5^j-2.
  \end{equation*}
\end{itemize}
These are the exact reset equations encoded by
\texttt{S6Audit.ResetHeadEq} and used in
\texttt{reset\_\allowbreak head\_\allowbreak predecessor}.

\subsection{First-block terminal algebra}

Let $g_{\mathrm{prev}}$ be the state before the incoming step and let
$e\ge2$ be its weight:
\begin{equation*}
  5g_{\mathrm{prev}}+1=2^e\,g.
\end{equation*}
The first-block terminal state $r$ is
\begin{equation*}
  r=\frac{5g_{\mathrm{prev}}+1}{2}=2^{e-1}g.
\end{equation*}
With $k_0=0$, the previous terminal becomes
\begin{equation*}
  s=2^{e-1}g+1,
\end{equation*}
and the reset predecessor becomes the candidate
\begin{equation*}
  x=2^{e-1}g+\delta5^{j-1}.
\end{equation*}
This is \texttt{first\_\allowbreak block\_\allowbreak terminal\_\allowbreak
eq} and \texttt{candidateX\_\allowbreak of\_\allowbreak reset\_\allowbreak
and\_\allowbreak terminal}.

\subsection{Block-head identities}

The block-head numerator has the exact identity
\begin{equation*}
  A_j+5^jq=2^L\left(B+\delta5^j\right),
\qquad 0<B<5^j,\quad A_j<5^j.
\end{equation*}
Since $A_j<5^j$, the quotient structure forces
\begin{equation*}
  q=m+\delta2^L,\qquad m<2^L.
\end{equation*}
This is \texttt{reset\_\allowbreak q0\_\allowbreak form}.

Under the reset equation, the same numerator has the block-head form
\begin{equation*}
  A_j+5^jq
    =2^{W_p}\left(5^{k_0+1}s-4+\delta5^j\right).
\end{equation*}
This is \texttt{block\_\allowbreak head\_\allowbreak identity\_\allowbreak
of\_\allowbreak reset}.

\subsection{The q0 interval}

The interval condition
\begin{equation*}
  2^{W_p}\le q<2^{W_j}
\end{equation*}
is equivalent to the block-head bounds
\begin{equation*}
  5^j\le2^t r_j,\qquad r_j<5^j,
\end{equation*}
by \texttt{q0\_\allowbreak interval\_\allowbreak iff\_\allowbreak
rj\_\allowbreak bound}.  The two forms are used in different layers:
the first belongs to the unified-core ledger, and the second belongs
to the reset-window predicate.

\subsection{Segment equations}

For a full-orbit segment of length $d$ from $g=\Orbit_7(j-1)$ to
$x=\Orbit_7(j-1+d)$ with weights $u_1,\dots,u_d$ and prefix weights
$W_i$, the exact segment equation is
\begin{equation*}
  2^{W+e-1}g+2^W\delta5^{j-1}
    =5^d\,g+\sum_{i=0}^{d-1}2^{W_i}5^{d-1-i},
\end{equation*}
where $W=W_d$ is the segment total weight and $e$ is the incoming
step weight.  The derivation is the word-orbit identity
\begin{equation*}
  2^Wx=5^dg+\wordA(w)
\end{equation*}
with the substitution $x=2^{e-1}g+\delta5^{j-1}$.  The small-$d$
closed forms above are exactly what the exclusions use.

\subsection{Residue rules for rise steps}

A rise step of weight $1$ satisfies
\begin{equation*}
  2y=5x+1,
\end{equation*}
so $y\equiv3\pmod5$.  A rise step of weight $2$ satisfies
\begin{equation*}
  4y=5x+1,
\end{equation*}
so $y\equiv4\pmod5$.  These are the only residue constraints used at
the block level; they are formalised in
\texttt{rise\_\allowbreak step\_\allowbreak next\_\allowbreak
mod\_\allowbreak five}.

\subsection{Corrected reachability predicate}

The corrected reachability predicate
$\ResetWindow(j,k_0,t,\delta,s)$ carries:
\begin{itemize}[leftmargin=*]
\item the previous-terminal equation $s5^{k_0}=r+1$;
\item the general-orbit reachability $\GeneralOrbit(r)$;
\item the full-orbit reachability $\FullOrbit(r_j)$;
\item the exact reset equation
  $\mathrm{ResetHeadEq}(s,j,k_0,t,\delta,r_j)$;
\item the size bound $s<5^{j-1-k_0}$, the oddness of $s$, and
  $5\nmid s$.
\end{itemize}
The predicate fixes the same $j,k_0,t$ in all fields.  A witness that
only fixes the difference $j-k_0-1$ is weaker and is not sufficient
for the corrected window bound.

\subsection{Arithmetic counterexample}

The unconstrained window bound is false.  The witness
\begin{equation*}
  j=36,\ k_0=0,\ t=2,\ \delta=3,\
  s=2^{83}-3\cdot5^{35}
\end{equation*}
satisfies
\begin{equation*}
  5s+3\cdot5^{36}
    =5(2^{83}-3\cdot5^{35})+3\cdot5^{36}
    =5\cdot2^{83},
\end{equation*}
so the valuation is $83$, while the unconstrained bound claims at most
$82$.  The witness satisfies the size, parity, and coprimality
conditions but fails the reachability fields of the corrected
predicate.  This is why the corrected bound is a projection of the
unified core rather than a standalone arithmetic estimate.

\subsection{Status}

\begin{center}
\small
\begin{tabular}{@{}ll@{}}
\toprule
Statement & Status \\
\midrule
small-block molecule formulas & Math: proven \\
reset equations & Lean: compiled \\
first-block terminal equation & Lean: compiled \\
candidate parameterisation & Lean: compiled \\
q0 interval equivalence & Lean: compiled \\
block-head identity of reset & Lean: compiled \\
corrected reachability definition & Lean: compiled \\
\bottomrule
\end{tabular}
\end{center}
